\chapter{A Irradiação do Divino}

% REVISÃO (ortografia): capítulo em grafia do Acordo de 1990 ("refletir", "reações", "objetos", "perspetiva", "atividades", "coleção", "diretamente"); uniformizar com a grafia pré-Acordo dos capítulos 00–19.
% REVISÃO (tratamento): alterna entre "tu" (ex.: "se tiveres", "Talvez tenhas") e "você" (ex.: "Quando vê uma mulher", "Não se preocupe"). Uniformizar.
O que é o divino?

Podemos considerar-nos criaturas puramente instintivas, porque temos um corpo
animal com a mesma natureza instintiva que um animal. A sobrevivência e a
procriação são tão fortes em nós como o são em gatos, cães e lobos. Mas também
existe o divino. Isto é algo a que nos elevamos ou para o qual nos voltamos;
porque não é instintivo, não será algo que encontraremos a menos que o
procuremos deliberadamente.

Para refletir sobre a divindade, temos os quatro \emph{brahma-vihāras}, as
belas qualidades altruístas que se podem manifestar através da forma humana
quando não há ego. Quando não estamos presos a comportamentos instintivos ou
reações emocionais baseadas na ignorância; quando há desapego e todo esse
processo de visão do eu cessa, então a divindade torna-se óbvia. Então, a
bondade, a compaixão, a alegria solidária e a serenidade da mente não são algo
que temos de alcançar, mas algo que se manifesta através destas formas.

Nas nossas vidas como seres separados, relacionamo-nos com as coisas. Como
seres individuais, temos relações com as coisas, temos de nos encontrar,
contactar e reagir ou responder a objetos constantemente durante o resto das
nossas vidas. No plano físico, temos de responder à presença uns dos outros de
alguma forma, seja ignorando, abraçando, prestando homenagem ou amaldiçoando. No
relacionamento, quando não há ego, então há essa divindade que se manifesta.
Assim, pode ver-se que a forma humana é uma forma para o divino.

Por outro lado, podemos pensar que é apenas para o nosso próprio bem: «É a minha
vida e posso fazer o que quiser, tenho direito à felicidade» -- e todo esse tipo
de egoísmo. Num nível instintivo, se não nos elevarmos acima da animalidade,
podemos simplesmente viver seguindo muito o instinto ou a emoção, ou podemos
viver num mundo de ideias e apego a ideias de como as coisas deveriam ser, o que
é um grande problema no mundo ocidental. Mas à medida que penetramos nisso e
vemos o sofrimento que advém de agarrar-se a qualquer coisa, à medida que essa
compreensão clara conduz ao desapego e o não apego, surge então uma resposta à
forma como as coisas são, que pode ser dividida nestas quatro categorias de
\emph{brahma-vihāras}.

\emph{Mettā, karuṇā, muditā e upekkhā} proporcionam uma reflexão: formam uma
sequência de como nos relacionarmos com o reino humano, com o reino animal e com
a natureza. \emph{Mettā} é, em grande medida, também a forma como nos devemos
relacionar connosco próprios. É a forma de nos relacionarmos connosco próprios
com bondade e aceitação, em vez de com aversão e julgamento. \emph{Mettā}
implica que aceitemos algo que pode não ser muito agradável -- como, por
exemplo, se tiveres dores físicas ou aspetos menos agradáveis do teu corpo ou do
teu caráter. Talvez tenhas muitas imperfeições, mau feitio ou algo do género. Se
tiveres \emph{mettā}, isso significa que consegues aceitar essas coisas tal como
são. Não as estás a julgar nem a condenar a partir de uma posição de ignorância
ou de complexos. Tens consciência de que são dolorosas, desagradáveis ou feias
-- mas a prática de \emph{mettā} é a capacidade de aceitar pacientemente as
falhas, a dor, as irritações e as frustrações dentro das nossas mentes e corpos,
bem como as coisas desagradáveis e irritantes que lhes incidem a partir do
exterior.

Isto porque, com \emph{mettā}, tais coisas já não são vistas como pessoais; não
há «eu» e «tu», não há «Tu fizeste-me isto\ldots{}» ou «Eu fiz-te isto».
\emph{Mettā} é ter perspetiva e não criar um problema, mesmo em relação à
injustiça, às inadequações e assim por diante, de nós próprios, dos outros ou da
sociedade. Isso não significa que não reparemos ou que não consigamos ver; mas
significa que não criamos problemas em torno disso, não carregamos isso nas
nossas mentes com amargura, ressentimento, raiva e tendências destrutivas. Com
\emph{mettā}, há sempre a capacidade de perdoar e recomeçar, de reconhecer as
coisas como elas são e não esperar que tudo se encaixe nos ideais que temos
sobre como as coisas deveriam ser. Isso não significa que nos resignemos
fatalisticamente à mediocridade, à tirania e à estupidez, mas significa que não
ficamos presos no padrão de ignorância que condiciona as formações mentais.
Assim, podemos suportar as vicissitudes da vida com bondade e aceitação.

Depois, há o \emph{karuṇā}. \emph{Karuṇā} é compaixão. Quando vemos o sofrimento dos
outros e a injustiça e a desigualdade que existem, respondemos com compaixão,
mas não é como uma pessoa rica a sentir pena dos pobres; não é isso. Não é
menosprezar os pobres, não é ser condescendente ou sentir pena das pessoas, mas
é compreender a situação difícil da nossa condição humana e tudo o que isso
implica. É a partir da compreensão da natureza do sofrimento, como surge e
cessa, que se pode ter verdadeira \emph{karuṇā} pelos outros seres.

Os britânicos têm muita \emph{karuṇā} pelos animais, não têm? A Grã-Bretanha é
um país bastante impressionante quando pensamos na quantidade de vida selvagem
que existe nesta área densamente povoada do sul de Inglaterra. Essa é uma boa
qualidade: \emph{karuṇā}. Portanto, a Grã-Bretanha é um país bondoso, onde as
pessoas, em geral, desenvolveram compaixão -- preocupação pelos infelizes e
pelos desfavorecidos.

Quando nos mudámos para Chithurst\footnote{Os residentes em West Sussex que se
  tornaram o primeiro mosteiro budista da floresta na Grã-Bretanha em 1979.},
havia pessoas que não nos queriam ali; mas a maioria da população local tentava
ser justa. Tinham compaixão por nós, por outras palavras, uma certa dose de
compaixão. Não nos iam fazer mal nem tentar livrar-se de nós, apesar de poderem
ter preferido que um belo mosteiro cristão ou uma família simpática e de classe
alta comprasse a Chithurst House, uma família que criasse cavalos e jogasse
pólo. Isso teria estado mais de acordo com o estado de espírito geral de West
Sussex, porque as pessoas gostam do que estão habituadas. Mas, como já havia
\emph{mettā-karuṇā} desenvolvida, houve apenas algumas pessoas que foram
diretamente hostis ou que tomariam alguma atitude contra nós. Portanto, pode-se
considerar isto como \emph{mettā-karuṇā}.

Às vezes, no budismo Theravāda, fica-se com a impressão de que não se deve
apreciar a beleza. Se vir uma flor bonita, deve-se contemplar a sua decadência;
ou se vir uma mulher bonita, deve-se contemplá-la como um cadáver em
decomposição. Isto tem um certo valor a um nível, mas não é uma posição fixa a
adotar. Não é que devamos sentir-nos compelidos a rejeitar a beleza e a
insistir na sua impermanência e em como ela se transforma em algo não tão belo
-- e depois francamente repulsivo. Essa é uma boa reflexão sobre \emph{anicca,
  dukkha e anatta}; mas pode deixar a impressão de que a beleza deve ser
refletida apenas em termos dessas três características, em vez de em termos da
experiência da beleza. Esta é a alegria de \emph{muditā} -- ser capaz de
apreciar a beleza nas coisas à nossa volta.

As flores são muito mais bonitas do que nós; admitimos que são mais bonitas;
esperamos que o sejam; não invejamos a sua beleza. Mas podemos realmente odiar
outra pessoa por ser bonita, porque isso é uma ameaça. A beleza de outra pessoa
\emph{faz-me} parecer menos bonita. Sou \emph{eu}, não é? Isto é ser observado;
não tentar forçar uma espécie de felicidade falsa à situação, mas deixar que
estas coisas cessem na tua mente. Estar claramente consciente deste problema
específico é permanecer com ele e não criar um problema a partir da visão do eu.
Depois, reconhece-o como \emph{anicca, dukkha, anatta}; deixa-o cessar;
deixa-o ir. Então, nesse desapego, encontramos uma alegria nos talentos, na
bondade e na beleza dos outros seres.

Quando olhas para as flores, experimentas um sentimento de alegria e isso é
muditā -- estás a regozijar-te ou a alegrar-te com a beleza de algo. Talvez
nunca tenhas refletido assim. Vemos coisas belas na natureza e, como não
representam uma ameaça para nós, podemos regozijar-nos com o pôr-do-sol ou com
a beleza das árvores, das montanhas e dos rios. Então, isso é \emph{muditā}: um
regozijo na beleza, na bondade e na verdade. E regozijamo-nos com a bondade dos
outros. Quando alguém faz algo de bom, ou ouvimos falar de alguma ação nobre, de
algum esforço heróico ou de algum sacrifício pessoal, surge uma sensação de
\emph{muditā}. Isso é alegria, alegria solidária.

Mas onde tendemos a falhar nisto é quando se torna uma questão de «tu» e «eu».
Podemos ficar com muita inveja da saúde e da beleza de alguém se estivermos
presos à visão do eu. Podemos sentir alegria pelas flores no jardim, mas depois
vamos à casa da vizinha e as flores dela são mais bonitas do que as nossas.
Podemos sentir inveja porque, a partir de uma posição egocêntrica, pensamos: «As
flores dela são mais bonitas do que as minhas, e ela é mais bonita do que eu»;
ou: «Ele é mais bonito do que eu, é mais inteligente»; ou: «Ele tem uma
personalidade melhor» -- tudo isto. Assim, sofremos com esta inveja e ciúme. É
um problema muito comum; na verdade, muitos seres humanos estão realmente presos
na inveja e no ciúme.

Se fôssemos à casa de uma pessoa rica, com os seus belos jardins, a piscina, os
belos tapetes orientais, o mobiliário encantador, a pessoa altruísta poderia
regozijar-se por estar num lugar bonito. Também se poderia pensar: «Hum, as
pessoas ricas provavelmente conseguiram isso enganando os pobres e explorando os
desfavorecidos\ldots{} resmungo, resmungo, resmungo.»

Lembro-me de uma vez ter entrado numa igreja com alguém em Londres e era uma
igreja linda. Eu disse: «Oh, que igreja encantadora.» Ele disse: «Sim,
provavelmente foi às custas de todas aquelas colónias que os britânicos
exploraram.» Eu não estava a comentar a história da igreja, mas sim a
experimentar verdadeiramente a alegria de estar num lugar bonito. E, no entanto,
podemos pensar que talvez essa igreja foi construída com os lucros do tráfico de
escravos ou do comércio de ópio. Talvez os traficantes de escravos e de drogas
do século passado se sentissem culpados, por isso construíram uma igreja
magnífica em Londres. No entanto, isso não significa que não seja bela, pois
não? Não a estamos a julgar do ponto de vista moral, mas a refletir sobre a
alegria, sobre a experiência da beleza, da bondade e da verdade: são estas que
trazem alegria às nossas vidas.

As pessoas que não conseguem ver a beleza do bem ou do verdadeiro são, no fundo,
amargas e mesquinhas, feias; vivem num reino feio onde não há regozijo na
beleza, na bondade e na verdade. Regozijar-se com estas coisas não significa
que nos deixemos levar por elas; a experiência da alegria deixa de ocorrer se
nos entregarmos à beleza e tentarmos agarrá-la, ou se nos agarrarmos à
experiência da alegria para tentar tê-la o tempo todo. Mas a \emph{muditā} é
certamente parte da nossa experiência humana.

\emph{Muditā} é a nossa capacidade de nos alegrarmos com a beleza e a
graciosidade das experiências da vida. É o sentimento de alegria, apreço e
gratidão pelas belezas e pelas coisas encantadoras da vida, e pelas coisas
encantadoras nas outras pessoas. Assim, quando não há ego, então há alegria;
encontramos alegria na bondade, na beleza das pessoas à nossa volta, na
sociedade ou nas condições naturais. Assim que tivermos discernimento,
descobrimos que desfrutamos e nos deleitamos com a beleza e a bondade das
coisas. A verdade, a beleza e a bondade encantam-nos; nelas encontramos
alegria: isso é \emph{muditā}.

Se vês a beleza como algo a agarrar, então ela desperta o desejo. Vês seres
humanos bonitos, uma mulher ou um homem bonito e pensas: «Eu quero-os.» Isso é
desejo. Isso não é regozijar-se com a beleza de alguém; é o desejo de possuir,
controlar e obter algo para ti próprio com isso.

Ao nível do instinto, é assim que as coisas são. É bastante natural. Se não nos
achássemos atraentes uns aos outros, ninguém iria querer procriar a espécie, não
é verdade? Se as atividades sexuais fossem dolorosas e miseráveis, ninguém iria
querer fazê-las. E se nos achássemos totalmente repulsivos e feios, então não
iríamos querer aproximar-nos uns dos outros, para não falar de algo tão íntimo
como o sexo. O desejo é o caminho natural nesse nível do reino sensorial. Não há
nada de errado nisso, mas existe a possibilidade de o ser humano o transcender.
Se o desejo fosse tudo o que somos e tudo o que podemos fazer, então deveríamos
segui-lo. Mas, como podemos transcender, temos esta ligação com o divino,
podemos elevar-nos acima da natureza grosseira e instintiva dos nossos próprios
corpos e do reino animal.

E é para isso que estou a apontar; não estou a condenar o reino animal. Os
animais podem trazer-nos muita alegria. Recentemente, em Chithurst, passei o
dia com a Doris, a nossa gata, e sempre senti que ela me trazia muita alegria. É
um animal muito agradável. Se me apego, no entanto, digo: «Tenho de ter a Doris.
Tenho de trazer a Doris para cá, para Amaravati. Não consigo viver sem ela.»
Depois, trago-a para cá, e ela teria de lutar com os gatos que vivem aqui, e
tudo isso seria apenas por minha causa, apenas para eu conseguir o que quero.
Nessa altura, já não seria uma experiência alegre. Traria muitos problemas.

Podemos refletir sobre como as coisas nos afetam. Querer sempre \emph{muditā} --
as belas flores, as cascatas e os belos pássaros a cantar -- significa que já
não podemos regozijar-nos com elas porque estamos a tentar agarrar-nos a elas.
Ficamos presos a todo o tipo de visões e opiniões sobre isso, por isso, mesmo
estando no meio disso, já não estamos realmente a desfrutar, a regozijar-nos
com isso, porque nos separámos disso através do nosso desejo por isso.

Na nossa vida como samaṇas\footnote{\emph{Samaṇa:} Significa «buscador
  espiritual»; neste contexto, alguém cuja prática o levou ao compromisso com
  uma forma de vida religiosa.}, contemplando a natureza, contemplando o Dhamma,
não temos de pensar que toda a beleza está lá apenas para nos corromper e nos
dar outro renascimento. Isso é outra visão do eu, não é? Mas esteja ciente de
como a beleza o afeta. Quando vê uma mulher bonita ou um homem bonito, como é
que isso afeta a sua mente? Pode haver uma atração inicial e, em seguida,
pode-se facilmente sentir-se ameaçado e rejeitá-la porque levamos uma vida de
celibato. Também pode dar uma segunda olhadela e perder-se nos pensamentos
sexuais que possam surgir desse contacto visual. No entanto, quanto mais
consciente estiver, menos tenderá a seguir as coisas como desejo, menos tenderá
a criar ou a alimentar os sentimentos com desejo e apego.

Portanto, a iluminação não significa uma espécie de indiferença insípida. Às
vezes, a iluminação é apresentada como se pudéssemos tornar-nos zombies sem
emoções, pessoas que já não sentem nada. Bem, enquanto houver um «eu», aquilo a
que chamamos alegria tende a estar tingida de egoísmo; fica manchada pelo nosso
«eu». Ficamos com inveja se temos algo bonito e alguém tem algo ainda mais
bonito, porque o egoísmo transforma sempre a beleza em possessividade, não é
verdade? Se as belezas da vida -- a alegria da verdade e a beleza da bondade --
provêm do eu, então estão sempre corrompidas pela inveja, pelo ciúme e por
pessoas invejosas.

Quando há egoísmo, mesmo que sejas o mais belo de todos, não é realmente uma
experiência alegre, porque estás sempre preocupado que possa haver alguém a
reivindicar essa coroa. Se adotares a visão do eu, há sempre essa possibilidade,
não é verdade? Mas quando não há eu, então a beleza não pertence a ninguém. Não
é minha nem tua; percebemos que não há possibilidade de a possuir de qualquer
forma, por isso não há desejo de possuir. Pode haver a alegria da experiência da
beleza sem que esta seja corrompida pelo egoísmo.

Depois, \emph{upekkhā}: equanimidade, serenidade. Para sermos capazes de
permanecer na serenidade da mente, não andamos por aí à procura de coisas belas
nas quais nos deleitarmos, porque não há eu. Respondes à beleza com alegria, mas
já não é algo que procuras ou que procuras como pessoa. A normalidade da vida é
\emph{upekkhā}, serenidade. Trata-se de ter paz com as dores e os sofrimentos
do processo de envelhecimento e com a separação dos entes queridos. Tudo isto é
a compreensão de \emph{upekkhā}, serenidade.

\emph{Upekkhā} não significa indiferença. Às vezes é traduzido como indiferença,
mas na verdade significa serenidade quando as coisas são feias, desagradáveis ou
comuns. Se seguir as práticas de \emph{asubha} -- observar, prestar atenção ao
que não é belo --, então começa a criar \emph{upekkhā}, equanimidade ou
serenidade.

Havia um hospital em Banguecoque que recebia todos os casos de homicídios e
mortes violentas, cadáveres encontrados nos canais e coisas do género. Se
entrasse lá numa segunda-feira, encontraria uma coleção do fim de semana e uma
variedade de objetos horríveis e macabros que, a princípio, lhe causariam uma
forte sensação de repulsa. Entrava e dizia: «Que nojo! Deixem-me sair daqui!»
-- porque, geralmente, não gostamos de olhar para corpos humanos que foram
mutilados e desfigurados e se encontram em estado de decomposição. A sociedade
civilizada mantém-se sempre afastada de tais coisas. Temos todas as
instituições para tratar disso, para que nunca tenha de chamar a nossa atenção.

Mas, na verdade, se meditarmos sobre isso, o resultado é equanimidade ou
serenidade. Se ultrapassarmos a aversão, o horror e a negatividade iniciais em
relação ao cadáver humano que está a apodrecer ou que foi cortado numa autópsia,
o resultado é equanimidade ou \emph{upekkhā}, uma tremenda paz e serenidade. Não
depressão. Não é aversão. Quando não há «eu», é possível permanecer num estado
de serenidade. Se houver «eu», então dizemos: «Detesto isto, não gosto, tira-me
isto, não aguento. Não consigo suportar isto. É repugnante, é nojento\ldots{}»
Mas quando há equanimidade, \emph{upekkhā}, não há eu. Assim, não se criam
problemas em relação ao processo de viver e à forma como as coisas se movem,
mudam e passam da beleza à decadência. Com \emph{muditā}, encontra-se alegria
na beleza e, quando a beleza desvanece, há equanimidade em vez de tristeza,
serenidade em vez de tristeza e desespero pela perda da beleza.

\emph{Upekkhā} é a capacidade de não seguir a aversão e de não se deixar levar
quando se vê coisas belas. Assim, não andamos apenas a correr de um lado para o
outro a tentar regozijar-nos com a beleza ou a sentir \emph{karuṇā} por todas
as criaturas infelizes. Podemos permitir a espera quando não há grande coisa a
acontecer. Com \emph{upekkhā}, não é preciso procurar algo que nos faça felizes,
uma causa pela qual lutar, ou envolver-se em atividades compulsivas. Este é
outro grande problema da humanidade moderna: tentamos canalizar a atividade
inquieta para boas causas, porque não há \emph{upekkhā}.

Tradicionalmente, os \emph{brahma-vihāras} são considerados \emph{lokiya
  dhamma,} Dhamma mundano, e não o Dhamma transcendente ou \emph{lokuttara}.
Devido à forma como a mente tende a pensar, surge a visão de que não vale a pena
preocupar-se com eles. «O \emph{lokuttara dhamma} é o importante. Não se dá
muita atenção aos \emph{lokiya dhammas}.»

Mas com a atenção plena, estamos na relação entre o \emph{lokiya} e o
\emph{lokuttara dhamma}. Relacionamo-nos ao nível do \emph{lokiya dhamma}
através dos \emph{brahma-vihāras} -- \emph{mettā, karuṇā, muditā} e
\emph{upekkhā}. Quando não há ego, quando não há ignorância a condicionar as
formações mentais, então existe a maneira como as coisas são -- os \emph{lokiya
  dhammas.} Mas não estamos a pedir que o \emph{muditā} seja uma experiência
permanente. Não esperamos ter uma experiência contínua, absoluta e eterna de
regozijo e alegria nas nossas vidas, porque não estamos apegados a isso como um
ponto de vista.

Assim, os \emph{brahma-vihāras} representam uma resposta espontânea a esta
experiência de nascimento e consciência quando não há um eu. São uma resposta
espontânea da abnegação, do \emph{anatta}, em vez de uma reação impulsiva do
desejo. Há uma diferença entre uma resposta espontânea à sabedoria e à atenção
plena e uma reação impulsiva ao desejo. A diferença reside nessa visão de um eu.
Na visão do eu, a pessoa continua a agarrar-se, apenas a reagir impulsivamente
com desejo às adversidades e experiências da vida. Quando já não há ignorância,
então há espontaneidade.

É isso que é a espontaneidade. Não há «eu» nela. É apenas uma forma cada vez
mais natural de responder à beleza, à verdade e à virtude; ou à dor e à miséria;
ou ao inverno, à primavera, ao verão e ao outono; aos seres afortunados ou aos
infelizes; e até mesmo à espera, segurando a tua chávena de chá, olhando pela
janela para a chuva.

Isto é apenas a contemplação do que é a divindade. Se refletires sobre a
natureza instintiva, o corpo ligado à terra, os seus desejos sexuais, as
capacidades reprodutivas, a sobrevivência, comer, beber, dormir, todas estas
% REVISÃO: "é apenas a forma como uma forma como esta sobrevive" — repetição de "forma como"; reformular, p. ex. "é apenas a maneira como uma forma como esta sobrevive".
necessidades instintivas básicas, não há nada de mal nisso; é apenas a forma
como uma forma como esta sobrevive. Tem de se reproduzir, não tem? Na verdade,
os seres humanos estão a tornar-se demasiado bons a procriar-se. É bastante
assustador, não é? Quantos milhares de milhões tem a população mundial? Quatro
ou cinco mil milhões neste planeta? E se fossem todos como animais, agindo
apenas por instinto, isso significaria quatro ou cinco mil milhões de seres
humanos egoístas, subdesenvolvidos, neuróticos e perturbados. Terrivelmente
assustador, não é? Podemos levar isso ao extremo oposto -- cinco mil milhões de
seres humanos iluminados! Ora, isso talvez não seja assim tão mau! Cinco mil
milhões de seres humanos iluminados em vez de cinco mil milhões de seres humanos
ignorantes e egoístas. Cinco mil milhões de seres humanos que conseguem
manifestar o divino nas suas vidas quotidianas, através de \emph{mettā, karuṇā,
  muditā, upekkhā.} Isso não soa assim tão mal, pois não? Soa bastante bem.

Mas cinco mil milhões de seres humanos a manifestar ganância, ódio e ilusão é um
quadro bastante sombrio. No entanto, não temos o direito de comentar sobre
\emph{eles}: este aqui, é o que temos, é nisto que podemos trabalhar. Não se
preocupe com os outros. Este aqui é o que pode realmente desenvolver através da
reflexão e da meditação.

\photoPage{image-015-festival-offerings.jpg}{Recebendo comida de esmola num dia de festival, Amaravati}{}
